\section{Semantics}
\label{sec:semantics}

\subsection{Expressions}
\label{sub:expr-sem}

\begin{definition}[expression evaluation]
\label{def:aeval}
Evaluation is defined by recursion on syntax, in two clauses.

\begin{enumerate}
% @interproof anchor def:aeval:arith
\item For the arithmetic expressions of \Cref{def:aexp},
\[
  % @interproof anchor def:aeval:arith:lit = "\sem{n}\sigma = n"
  \sem{n}\sigma = n, \qquad
  % @interproof anchor def:aeval:arith:var = "\sem{x}\sigma = \sigma(x)"
  \sem{x}\sigma = \sigma(x), \qquad
  % @interproof anchor def:aeval:arith:add = "\sem{a_1 + a_2}\sigma = \sem{a_1}\sigma + \sem{a_2}\sigma"
  \sem{a_1 + a_2}\sigma = \sem{a_1}\sigma + \sem{a_2}\sigma .
\]
% @interproof anchor def:aeval:bool
\item For the boolean expressions of \Cref{def:bexp}, writing
$\mathbb{B} = \{\mathsf{tt}, \mathsf{ff}\}$,
\begin{align*}
% @interproof anchor def:aeval:bool:tt = "\sem{\kw{true}}\sigma &= \mathsf{tt}"
  \sem{\kw{true}}\sigma &= \mathsf{tt}, &
% @interproof anchor def:aeval:bool:le = "\sem{a_1 \le a_2}\sigma &= (\sem{a_1}\sigma \le \sem{a_2}\sigma)"
  \sem{a_1 \le a_2}\sigma &= (\sem{a_1}\sigma \le \sem{a_2}\sigma), \\
% @interproof anchor def:aeval:bool:not = "\sem{\lnot b}\sigma &= \lnot \sem{b}\sigma"
  \sem{\lnot b}\sigma &= \lnot \sem{b}\sigma, &
% @interproof anchor def:aeval:bool:and = "\sem{b_1 \land b_2}\sigma &= \sem{b_1}\sigma \land \sem{b_2}\sigma"
  \sem{b_1 \land b_2}\sigma &= \sem{b_1}\sigma \land \sem{b_2}\sigma .
\end{align*}
\end{enumerate}
Both clauses define total functions: expressions have no side effects and
cannot fail.
\end{definition}

\subsection{Commands}
\label{sub:cmd-sem}

\begin{definition}[big-step evaluation]
\label{def:bigstep}
The relation $\bigstep{c}{\sigma}{\tau}$ --- ``running $c$ in $\sigma$
terminates in $\tau$'' --- is inductively generated by
\[
% @interproof anchor def:bigstep:skip = "\frac{}{\bigstep{\skipk}{\sigma}{\sigma}}"
  \frac{}{\bigstep{\skipk}{\sigma}{\sigma}}
  \qquad
% @interproof anchor def:bigstep:assign = "\frac{}{\bigstep{x := a}{\sigma}{\upd{\sigma}{x}{\sem{a}\sigma}}}"
  \frac{}{\bigstep{x := a}{\sigma}{\upd{\sigma}{x}{\sem{a}\sigma}}}
  \qquad
% @interproof anchor def:bigstep:seq = "\frac{\bigstep{c_1}{\sigma}{\tau} \quad \bigstep{c_2}{\tau}{\upsilon}} {\bigstep{c_1 ; c_2}{\sigma}{\upsilon}}"
  \frac{\bigstep{c_1}{\sigma}{\tau} \quad \bigstep{c_2}{\tau}{\upsilon}}
       {\bigstep{c_1 ; c_2}{\sigma}{\upsilon}}
\]
\[
% @interproof anchor def:bigstep:iteTrue = "\frac{\sem{b}\sigma = \mathsf{tt} \quad \bigstep{c_1}{\sigma}{\tau}} {\bigstep{\ifk\ b\ \thenk\ c_1\ \elsek\ c_2}{\sigma}{\tau}}"
  \frac{\sem{b}\sigma = \mathsf{tt} \quad \bigstep{c_1}{\sigma}{\tau}}
       {\bigstep{\ifk\ b\ \thenk\ c_1\ \elsek\ c_2}{\sigma}{\tau}}
  \qquad
% @interproof anchor def:bigstep:iteFalse = "\frac{\sem{b}\sigma = \mathsf{ff} \quad \bigstep{c_2}{\sigma}{\tau}} {\bigstep{\ifk\ b\ \thenk\ c_1\ \elsek\ c_2}{\sigma}{\tau}}"
  \frac{\sem{b}\sigma = \mathsf{ff} \quad \bigstep{c_2}{\sigma}{\tau}}
       {\bigstep{\ifk\ b\ \thenk\ c_1\ \elsek\ c_2}{\sigma}{\tau}}
\]
\[
% @interproof anchor def:bigstep:whileFalse = "\frac{\sem{b}\sigma = \mathsf{ff}} {\bigstep{\whilek\ b\ \dok\ c}{\sigma}{\sigma}}"
  \frac{\sem{b}\sigma = \mathsf{ff}}
       {\bigstep{\whilek\ b\ \dok\ c}{\sigma}{\sigma}}
  \qquad
% @interproof anchor def:bigstep:whileTrue = "\frac{\sem{b}\sigma = \mathsf{tt} \quad \bigstep{c}{\sigma}{\tau} \quad \bigstep{\whilek\ b\ \dok\ c}{\tau}{\upsilon}} {\bigstep{\whilek\ b\ \dok\ c}{\sigma}{\upsilon}}"
  \frac{\sem{b}\sigma = \mathsf{tt} \quad
        \bigstep{c}{\sigma}{\tau} \quad
        \bigstep{\whilek\ b\ \dok\ c}{\tau}{\upsilon}}
       {\bigstep{\whilek\ b\ \dok\ c}{\sigma}{\upsilon}}
\]
using the update notation of the companion note and the evaluation of
\Cref{def:aeval}. A non-terminating run is simply the absence of a derivation.
\end{definition}

\begin{lemma}[determinism]
\label{lem:bigstep-det}
If $\bigstep{c}{\sigma}{\tau}$ and $\bigstep{c}{\sigma}{\tau'}$ then
$\tau = \tau'$.
\end{lemma}

\begin{proof}
Induction on the first derivation, inverting the second at each step. Every
pair of rules of \Cref{def:bigstep} with the same conclusion is separated by
the value of $\sem{b}\sigma$, so at most one applies.
\end{proof}
